\documentclass[12pt,a4paper]{article}
\usepackage[utf8]{inputenc}
\usepackage[russian]{babel}
\usepackage[T2A]{fontenc}
\usepackage{amsmath, amssymb, amsfonts}
\usepackage{graphicx}
\usepackage{booktabs}
\usepackage{hyperref}
\usepackage{geometry}
\geometry{margin=2.5cm}
\usepackage{caption}
\usepackage{float}
\usepackage{siunitx}

\title{Decoupled Observer Patch Holography (dOPH): \\
Математическое обоснование и первые астрофизические тесты}
\author{v04226079}
\date{\today}

\begin{document}

\maketitle

\begin{abstract}
Представлена модель \textbf{Decoupled Observer Patch Holography (dOPH)} — scale-dependent модифицированная теория гравитации, основанная на голографических принципах и динамическом decoupling наблюдательных патчей.

Ключевой результат работы — вывод универсального scaling factor 
\[
P = \sqrt{\frac{8}{3}} \approx 1.633
\]
как геометрического fixed point из минимизации функционала несоответствия на интерфейсах патчей.

Модель протестирована на массивной эллиптической галактике NGC 4649 (M60). При реалистичном отношении масса/светимость звёзд (\( M/L \approx 8 \)) dOPH успешно воспроизводит практически плоский профиль дисперсии скоростей глобулярных кластеров до \( \sim 38 \) кпк. Теория показывает конкурентные результаты по сравнению с MOND и моделями с тёмной материей, используя существенно меньше свободных параметров.
\end{abstract}

\section{Введение}

Наблюдаемая динамика галактик и скоплений остаётся серьёзным вызовом для современной физики. В данной работе развивается框架 Decoupled Observer Patch Holography (dOPH).

\section{Математическое обоснование \( P = \sqrt{8/3} \)}

На интерфейсе двух соседних observer patches существуют естественные ускорения \( g_1 \) и \( g_2 \). Минимизируется функционал:

\begin{equation}
J(\lambda) = (g_1 - \lambda g_2)^2 + \beta \left( \lambda g_2 - \frac{g_1}{\lambda} \right)^2
\end{equation}

При \( g_1 \approx g_2 = g \) и \( \beta = 3/8 \) минимум достигается точно при \( \lambda = P = \sqrt{8/3} \).

Данное значение является **геометрическим fixed point** теории, вытекающим из оптимальной 3D-упаковки (HCP/FCC), минимизации mismatch и конформных соображений.

\section{Тест на NGC 4649 (M60)}

Галактика NGC 4649 демонстрирует почти плоский профиль дисперсии глобулярных кластеров (\( \sigma_{\rm LOS} \approx 220 \)--$235$ км/с до \( \sim 38 \) кпк). Модель dOPH при \( M/L \approx 8 \) успешно воспроизводит наблюдаемые данные, значительно превосходя чистую ньютоновскую модель со звёздной массой.

\section{Обсуждение и дальнейшие планы}

Модель хорошо работает на галактических масштабах, но требует доработки для описания столкновений скоплений (Bullet Cluster). В будущем планируется развитие non-local членов, multi-component decoupling и временной зависимости.

\section{Заключение}

Получено строгое математическое обоснование \( P = \sqrt{8/3} \) как геометрического fixed point. Тест на NGC 4649 показал перспективность подхода. Модель представляет собой интересную голографическую scale-dependent альтернативу стандартной тёмной материи.

Весь код открыт: \url{https://github.com/v04226079-sketch/dOPH-Decoupled-Observer-Patch-Holography}

\end{document}
